\esSezione{Rete di interconnessione 16:4 su board}\label{sec:esercizio1-3}

\begin{traccia}
Sintetizzare ed implementare su board Nexys A7 la rete di interconnessione 16:4 sviluppata al punto precedente. Le selezioni sono fornite tramite switch e le 4 uscite sono visualizzate sui led. Il dato di 16 bit è immesso 8 bit alla volta con gli switch, mediante una rete di acquisizione che sfrutta due pulsanti della board per caricare, rispettivamente, la prima e la seconda metà del dato.
\end{traccia}

\progetto

Alla rete di interconnessione, già presentata nell'esercizio~1.2 (listato~\ref{lst:esercizio1-2:interconn16-4}) e costruita con il multiplexer 16:1 (listato~\ref{lst:esercizio1-1:mux-16-1}) e il demultiplexer 1:4 (listato~\ref{lst:esercizio1-2:demux-1-4}), si aggiunge un registro di acquisizione, \texttt{acq\_reg}. Il registro memorizza il dato a 16 bit da inviare agli ingressi \texttt{x} della rete, in modo che esso rimanga disponibile dopo che gli switch sono stati riposizionati per inserire l'altra metà. Il top module \texttt{top\_1\_3} collega i tre elementi e mette in corrispondenza le porte con le risorse della board.

L'assegnazione degli switch è la seguente:
\begin{enumerate}
  \item \texttt{SW(7 downto 0)}: otto bit del dato, che costituiscono la metà da caricare;
  \item \texttt{SW(11 downto 8)}: selezione dell'ingresso, \texttt{s\_i};
  \item \texttt{SW(13 downto 12)}: selezione dell'uscita, \texttt{s\_o}.
\end{enumerate}
Il pulsante \texttt{BTNL} carica la prima metà del dato (bit da 0 a 7), il pulsante \texttt{BTNR} la seconda (bit da 8 a 15). Le quattro uscite della rete sono mostrate sui led \texttt{LED(3 downto 0)}.

Il multiplexer e il demultiplexer sono utilizzati senza modifiche e non vengono ripresentati.

\implementazione

L'entità \texttt{acq\_reg} ha come ingressi il clock \texttt{clk}, i due comandi di caricamento \texttt{load\_lo} e \texttt{load\_hi} e il vettore \texttt{d} di 8 bit; l'uscita \texttt{q} è un vettore di 16 bit. L'architettura è comportamentale e contiene un solo process con sensitivity list limitata a \texttt{clk}: sul fronte di salita, se \texttt{load\_lo} vale \texttt{'1'}, i bit da 0 a 7 del segnale interno \texttt{r} assumono il valore di \texttt{d}; analogamente, se \texttt{load\_hi} vale \texttt{'1'}, i bit da 8 a 15 di \texttt{r} ricevono \texttt{d}. Il registro è inizializzato a zero e l'uscita \texttt{q} riporta in modo continuo il contenuto di \texttt{r}. Le due metà sono indipendenti: il caricamento di una non altera l'altra.

\vhdlfile{esercizio1-3/acq_reg.vhd}{Implementazione registro di acquisizione}{lst:esercizio1-3:acq-reg}

Il top module ha come porte il clock a 100 MHz \texttt{CLK100MHZ}, gli switch \texttt{SW} a 14 bit, i pulsanti \texttt{BTNL} e \texttt{BTNR} e i led \texttt{LED} a 4 bit. L'architettura è strutturale. Due cicli \texttt{for generate} collegano gli otto switch meno significativi al segnale \texttt{d\_half} e il segnale \texttt{y} ai led. Il registro \texttt{acq\_reg} è istanziato con \texttt{BTNL} e \texttt{BTNR} come comandi di caricamento e fornisce il segnale \texttt{data}; la rete \texttt{interconn16\_4} riceve \texttt{data} sull'ingresso \texttt{x}, \texttt{SW(11 downto 8)} su \texttt{s\_i}, \texttt{SW(13 downto 12)} su \texttt{s\_o} e produce \texttt{y}.

\vhdlfile{esercizio1-3/top_1_3.vhd}{Implementazione su Nexys A7 della rete di interconnessione 16:4}{lst:esercizio1-3:top}

\sintesiboard

L'uso della board prevede due fasi. Dapprima si impostano gli otto switch \texttt{SW(7 downto 0)} sulla prima metà del dato e si preme \texttt{BTNL}; poi si impostano gli stessi switch sulla seconda metà e si preme \texttt{BTNR}. Il dato a 16 bit risulta così memorizzato nel registro. Con \texttt{SW(11 downto 8)} si sceglie l'ingresso da propagare e con \texttt{SW(13 downto 12)} l'uscita su cui renderlo disponibile: il bit scelto compare sul corrispondente led tra \texttt{LED(0)} e \texttt{LED(3)}, mentre gli altri restano a zero. Gli indici di selezione sono espressi in notazione posizionale, come nelle simulazioni dell'esercizio~1.2.

Nel file dei constraint sono state decommentate le seguenti righe, relative al clock, agli switch da \texttt{SW[0]} a \texttt{SW[13]}, ai quattro led e ai due pulsanti.

\xdcfile{esercizio1-3/vincoli.xdc}{Constraint}{lst:esercizio1-3:vincoli}
